\documentclass[9pt,letterpaper]{article}
\usepackage[margin=0.2in]{geometry}
\usepackage{amsmath,amssymb}
\usepackage{enumitem}
\usepackage{booktabs}
\usepackage{multicol}
\usepackage{titlesec}
\usepackage{parskip}
\usepackage{tcolorbox}
\usepackage{array}

\titlespacing*{\section}{0pt}{4pt}{1pt}
\titlespacing*{\subsection}{0pt}{3pt}{1pt}
\titleformat{\section}{\normalsize\bfseries}{}{0em}{}
\titleformat{\subsection}{\small\bfseries}{}{0em}{}

\setlist[itemize]{noitemsep, topsep=1pt, leftmargin=*, parsep=0pt}
\setlist[enumerate]{noitemsep, topsep=1pt, leftmargin=*, parsep=0pt}

\newtcolorbox{calcbox}{colback=gray!8, colframe=gray!50, boxrule=0.3pt, left=2pt, right=2pt, top=1pt, bottom=1pt, before skip=2pt, after skip=2pt, fontupper=\small}
\newtcolorbox{keybox}{colback=blue!5, colframe=blue!40, boxrule=0.3pt, left=2pt, right=2pt, top=1pt, bottom=1pt, before skip=2pt, after skip=2pt, fontupper=\small}

\pagestyle{plain}
\setlength{\parindent}{0pt}
\setlength{\parskip}{1pt plus 0.5pt minus 0.5pt}
\setlength{\abovedisplayskip}{3pt}
\setlength{\belowdisplayskip}{3pt}
\small

\begin{document}

\vspace*{-0.25in}
\begin{center}
{\large \textbf{Midterm 1 Comprehensive Study Guide --- Stat 490}}\\[-2pt]
{\scriptsize Design of Experiments $\cdot$ Spring 2026}
\end{center}

{\footnotesize 0. Table of contents: 1) Foundations\quad 2) One-way ANOVA\quad 3) Multiple Comparisons\quad 4) Model Adequacy\quad 5) Random Effects\quad 6) RCBD\quad 7) Power\quad 8) Non-parametrics\quad 9) Design Principles\quad 10) Distributions}
\vspace{2pt}\hrule\vspace{2pt}

%% ============================================================
\section{1. Foundations: Hypothesis Testing \& Two-Sample Inference}
%% ============================================================

\subsection{1.1 Core Definitions}
\begin{itemize}
    \item Population: full set of units. Sample: subset used for inference. SRS: every unit has equal chance.
    \item $H_0$ (null): no effect/difference. $H_a$ (alternative): effect/difference exists.
    \item Type I error ($\alpha$): reject $H_0$ when true (false positive). Type II ($\beta$): fail to reject $H_0$ when $H_a$ true (false negative).
    \item p-value: $P(\text{test stat as or more extreme} \mid H_0 \text{ true})$. Reject $H_0$ if p-value $< \alpha$.
    \item Power $=1-\beta = P(\text{reject } H_0 \mid H_a \text{ true})$. Increases with larger $n$, larger effect size, larger $\alpha$, smaller $\sigma$.
\end{itemize}

\subsection{1.2 $t$-Tests: Assumptions \& Theory}

\textbf{One-sample $t$}: test $H_0:\mu=\mu_0$.
\begin{itemize}
    \item Assumptions: SRS from population; errors $X-\mu$ i.i.d; population approx normal or $n$ large.
    \item Test statistic: $t=(\bar{X}-\mu_0)/(s/\sqrt{n})$; compares observed mean shift to its SE under $H_0$.
\end{itemize}

\textbf{Two-sample pooled $t$}: $H_0:\mu_1=\mu_2$ with equal variances.
\begin{itemize}
    \item Assumptions: two independent SRSs; both populations approx normal; $\sigma_1^2=\sigma_2^2$ (check SD ratio in $[1/2,2]$, F/Bartlett/Levene).
    \item Test stat: $t=(\bar{X}_1-\bar{X}_2)/(s_p\sqrt{1/n_1+1/n_2})$, where $s_p^2$ pools both sample variances.
\end{itemize}

\textbf{Welch (non-pooled) $t$}: $H_0:\mu_1=\mu_2$ with unequal variances.
\begin{itemize}
    \item Assumptions: two independent SRSs; normal or large $n$; variance may differ but shapes similar.
    \item Test stat: $t=(\bar{X}_1-\bar{X}_2)/\sqrt{s_1^2/n_1+s_2^2/n_2}$ with Satterthwaite df.
\end{itemize}

\textbf{Paired $t$}: $H_0:\mu_d=0$ for differences.
\begin{itemize}
    \item Assumptions: pairs form SRS of differences; differences $d_i$ approx normal or $n$ large; pairs independent.
    \item Test stat: $t=\bar{d}/(s_d/\sqrt{n_d})$; reduces variability by blocking at subject level.
\end{itemize}

\begin{center}\small
\begin{tabular}{@{}llc@{}}
\toprule
Scenario & Test Statistic & df \\
\midrule
One-sample & $t=(\bar{x}-\mu_0)/(s/\sqrt{n})$ & $n-1$ \\
Pooled two-sample & $t=(\bar{x}_1-\bar{x}_2)/(s_p\sqrt{1/n_1+1/n_2})$ & $n_1+n_2-2$ \\
Welch two-sample & $t=(\bar{x}_1-\bar{x}_2)/\sqrt{s_1^2/n_1+s_2^2/n_2}$ & Satterthwaite \\
Paired & $t=\bar{d}/(s_d/\sqrt{n_d})$ & $n_d-1$ \\
\bottomrule
\end{tabular}
\end{center}

\begin{calcbox}
TI-84: \texttt{STAT $\to$ TESTS $\to$ 2:T-Test} (one-sample), \texttt{4:2-SampTTest} (two-sample). For pooled, set \texttt{Pooled:Yes}. Use \texttt{Stats} input to enter $\bar{x},s,n$ directly.
\end{calcbox}

\subsection{1.3 CIs \& Assumption Checks}

CI for $\mu$: $\bar{x} \pm t_{\alpha/2,df}\, s/\sqrt{n}$. CI for $\mu_1-\mu_2$: $(\bar{x}_1-\bar{x}_2) \pm t_{\alpha/2,df}\,SE$; if CI excludes 0 $\Rightarrow$ reject $H_0$.

\begin{itemize}
    \item Normality: histogram, QQ-plot, Shapiro--Wilk test. t/F tests robust to mild non-normality for balanced moderate $n$.
    \item Equal variance (2 groups): F-test on ratio $s_1^2/s_2^2$ (sensitive to non-normality).
    \item Equal variance ($k$ groups): Bartlett (powerful if normal, sensitive otherwise); Levene (robust, uses absolute deviations from median).
\end{itemize}

%% ============================================================
\section{2. One-Way ANOVA (Fixed Effects, CRD)}
%% ============================================================

\subsection{2.1 Model \& Assumptions}

One-factor fixed-effects ANOVA model (balanced CRD):
\[
y_{ij}=\mu+\tau_i+\varepsilon_{ij},\quad i=1,\ldots,a,\; j=1,\ldots,n,\quad \varepsilon_{ij}\sim\text{NID}(0,\sigma^2).
\]

Assumptions:
\begin{itemize}
    \item Independence: ensured by random assignment (no time/space patterns).
    \item Normality: error terms $\varepsilon_{ij}$ are normal within each group.
    \item Constant variance: $\text{Var}(\varepsilon_{ij})=\sigma^2$ for all $i,j$.
\end{itemize}

Treatment means: $\mu_i=\mu+\tau_i$, with constraint $\sum_{i=1}^a \tau_i=0$. Total $N=an$.

\subsection{2.2 Decomposition, EMS, \& ANOVA Table}

Key identity: $SS_T = SS_{\text{Trt}} + SS_E$.
\[
SS_T=\sum_{i,j}(y_{ij}-\bar{y}_{..})^2,\quad SS_{\text{Trt}}=\sum_i n(\bar{y}_{i.}-\bar{y}_{..})^2,\quad SS_E=\sum_{i,j}(y_{ij}-\bar{y}_{i.})^2.
\]

ANOVA table:
\begin{center}\small
\begin{tabular}{@{}lcccc@{}}
\toprule
Source & df & SS & MS & $F$ \\
\midrule
Treatment & $a-1$ & $SS_{\text{Trt}}$ & $MS_{\text{Trt}}=SS_{\text{Trt}}/(a-1)$ & $MS_{\text{Trt}}/MS_E$ \\
Error & $N-a$ & $SS_E$ & $MS_E=SS_E/(N-a)$ & \\
Total & $N-1$ & $SS_T$ & & \\
\bottomrule
\end{tabular}
\end{center}

Expected mean squares (fixed effects):
\[
E[MS_E]=\sigma^2,\qquad E[MS_{\text{Trt}}]=\sigma^2+\frac{n\sum_{i=1}^a \tau_i^2}{a-1}.
\]
Under $H_0$ (all $\tau_i=0$): both MS estimate $\sigma^2$, so $F\approx1$. Under $H_a$: $E[MS_{\text{Trt}}]>E[MS_E]$, so $F$ tends to be large.

Estimates: $\hat{\mu}=\bar{y}_{..}$, $\hat{\tau}_i=\bar{y}_{i.}-\bar{y}_{..}$, $\hat{\mu}_i=\bar{y}_{i.}$, residuals $e_{ij}=y_{ij}-\bar{y}_{i.}$. Coefficient of determination: $R^2=SS_{\text{Trt}}/SS_T$.

\subsection{2.3 Regression Equivalence}

ANOVA with $a$ levels is equivalent to regression with $a-1$ dummy variables. With reference level 1:
\[
y_{ij}=\beta_0+\beta_2x_{2,ij}+\cdots+\beta_a x_{a,ij}+\varepsilon_{ij},
\]
where $\beta_0$ is mean of level 1 and $\beta_i$ are differences from level 1. ANOVA $F$ for factor $=$ regression $F$ for all dummies.

\subsection{2.4 CIs from ANOVA}

Using $MS_E$ from ANOVA and error df $=N-a$:
\begin{itemize}
    \item CI for $\mu_i$: $\bar{y}_{i.} \pm t_{\alpha/2,N-a}\sqrt{MS_E/n}$.
    \item CI for $\mu_i-\mu_j$: $(\bar{y}_{i.}-\bar{y}_{j.}) \pm t_{\alpha/2,N-a}\sqrt{2MS_E/n}$ (balanced).
\end{itemize}

\subsection{2.5 Worked Example: Dog Food Shelf Height}

3 shelf heights (Knee, Eye, Waist), $n=8$ stores per group, $N=24$ total.

\begin{center}\small
\begin{tabular}{@{}lccc@{}}
\toprule
 & Knee & Eye & Waist \\
\midrule
$\bar{y}_{i.}$ & 81.0 & 84.625 & 90.875 \\
\bottomrule
\end{tabular}
\end{center}

$\bar{y}_{..}=85.5$, $SS_{\text{Trt}}=399.2$, $SS_E=288.7$, $MS_E=13.75$. Then $F=199.62/13.75=14.52$ with df $(2,21)$, p-value $\approx0.00011$; strong evidence of unequal means.

95\% CI for $\mu_{\text{Knee}}$: $81 \pm 2.08\sqrt{13.75/8}=81\pm2.73=(78.27,83.73)$.

\begin{calcbox}
TI-84: \texttt{Fcdf(14.52,1E99,2,21)} gives p-value. \texttt{invT(0.025,21)} gives $t_{0.025,21}$. Use \texttt{1-Var Stats} on each group list to get $\bar{x},s^2,n$.
\end{calcbox}

%% ============================================================
\section{3. Multiple Comparisons}
%% ============================================================

\subsection{3.1 Familywise Error \& Concepts}

With $K$ tests at level $\alpha$, FW (experiment-wise) Type I error $\approx1-(1-\alpha)^K$. For $a$ means, $K=\binom{a}{2}$ pairwise tests. For $a=4$ and $\alpha=0.05$, FW error $\approx0.26$.

Family definition matters: could be ``all pairwise comparisons'', a set of planned contrasts, or all contrasts considered after seeing data.

\subsection{3.2 Post-Hoc Procedures: Assumptions \& Use}

\begin{center}\small
\begin{tabular}{@{}p{1.7cm}p{4.8cm}p{3.6cm}c@{}}
\toprule
Method & Key Idea / Formula & When to Use & FWER? \\
\midrule
Fisher LSD & $\text{LSD}=t_{\alpha/2,df_E}\sqrt{2MS_E/n}$; pair $i,j$ sig if $|\bar{y}_{i.}-\bar{y}_{j.}|>\text{LSD}$. & After sig ANOVA $F$; few groups; liberal. & No \\
Bonferroni & Each test at $\alpha/K$ or CIs at $1-\alpha/K$. & Any small family of planned tests. & Yes \\
Tukey HSD & $\text{HSD}=q_{\alpha,a,df_E}\sqrt{MS_E/n}$. & All pairwise comparisons; equal variances; near-balanced. & Yes \\
Dunnett & Critical values for trt vs control differences only. & All trts vs single control. & Yes \\
Scheff\'e & Reject if $|\hat{\Gamma}|/SE>\sqrt{(a-1)F_{\alpha,a-1,df_E}}$. & Any (even post-hoc) contrast. & Yes \\
\bottomrule
\end{tabular}
\end{center}

Assumptions: all inherit ANOVA assumptions (normal, equal variance, independence). Tukey/Dunnett additionally assume equal $MS_E$ across groups; Scheff\'e valid for arbitrary contrasts, even selected after seeing data.

\subsection{3.3 Dog Food Pairwise Comparisons}

From \S2.5: $a=3$, $n=8$, $MS_E=13.75$, $df_E=21$. $SE=\sqrt{2MS_E/n}=1.854$.

\begin{center}\small
\begin{tabular}{@{}lcccc@{}}
\toprule
Pair & $|\bar{y}_{i.}-\bar{y}_{j.}|$ & LSD & Tukey HSD & Bonf. thresh \\
\midrule
W--K & 9.875 & 3.86 & 4.67 & 4.86 \\
W--E & 6.250 & 3.86 & 4.67 & 4.86 \\
K--E & 3.625 & 3.86 & 4.67 & 4.86 \\
\bottomrule
\end{tabular}
\end{center}

All methods find W differs from K and E; K vs E not significant.

\subsection{3.4 Contrasts \& Orthogonality}

A contrast is $\Gamma=\sum_{i=1}^a c_i\mu_i$ with $\sum c_i=0$. Estimate $\hat{\Gamma}=\sum c_i\bar{y}_{i.}$; $SE(\hat{\Gamma})=\sqrt{MS_E\sum c_i^2/n}$ (balanced). CI: $\hat{\Gamma}\pm t_{\alpha/2,df_E}SE(\hat{\Gamma})$.

Example: control vs average of two treatments: $C=(-2,1,1)$; protein effect in $2\times2$: $C=(0.5,0.5,-0.5,-0.5)$.

Two contrasts $C,D$ are orthogonal if $\sum c_id_i=0$ (balanced). There are at most $a-1$ mutually orthogonal contrasts; their SS add to $SS_{\text{Trt}}$ (each 1 df).

\subsection{3.5 Data Snooping \& Scheff\'e}

Choosing contrasts after looking at data inflates Type I error (implicit multiple testing). Scheff\'e's method controls FW error for \emph{all} contrasts at once; more conservative than Tukey for pairwise comparisons but valid under data snooping.

%% ============================================================
\section{4. Model Adequacy \& Transformations}
%% ============================================================

\subsection{4.1 Residual Diagnostics}

Residuals $e_{ij}=y_{ij}-\bar{y}_{i.}$ summarize unexplained variation. Checks:
\begin{itemize}
    \item Normality: QQ-plot; Shapiro--Wilk test.
    \item Constant variance: residuals vs fitted; look for funnels. Bartlett/Levene tests.
    \item Outliers: studentized residual $|e^*|>3$ suspicious.
    \item Independence: residuals vs run order for time trends.
\end{itemize}

\subsection{4.2 Variance-Stabilizing Transformations}

If $\text{Var}(Y)=g(\mu)$, find $h$ with $h'(\mu)\propto1/\sqrt{g(\mu)}$ so $\text{Var}(h(Y))$ is roughly constant (delta method). Common cases:

\begin{center}\small
\begin{tabular}{@{}llll@{}}
\toprule
Distribution & $\text{Var}(Y)$ & Transform & Example \\
\midrule
Poisson & $\mu$ & $h(y)=\sqrt{y}$ & Flaw counts \\
Binomial & $n\mu(1-\mu)$ & $h(y)=\sin^{-1}\sqrt{y}$ & Proportions \\
Lognormal-like & $\propto\mu^2$ & $h(y)=\ln y$ & Times, prices \\
\bottomrule
\end{tabular}
\end{center}

Decision rule: if variance grows with mean, try log or square root; if data are proportions near 0/1, try arcsin--sqrt. If transformations fail and assumptions still poor, use non-parametric tests.

\begin{calcbox}
TI-84: store raw data in \texttt{L1}, set \texttt{L2 = $\sqrt{}$(L1)} or \texttt{L2 = ln(L1)}, then run \texttt{1-Var Stats L2}.
\end{calcbox}

%% ============================================================
\section{5. Random Effects Model}
%% ============================================================

\subsection{5.1 Fixed vs Random Effects}

Fixed effects: treatment levels are the \emph{only} ones of interest; inference about those specific means. Random effects: levels are randomly sampled; inference about variance among levels and the population of levels.

Examples: specific 3 labs vs 3 labs sampled from all labs; specific therapists vs random sample of therapists.

\subsection{5.2 Model, Assumptions, EMS}

Random one-way ANOVA:
\[
y_{ij}=\mu+\tau_i+\varepsilon_{ij},\quad \tau_i\sim N(0,\sigma_\tau^2),\; \varepsilon_{ij}\sim N(0,\sigma^2),\; \text{independent}.
\]

Assumptions:
\begin{itemize}
    \item Levels $i$ are random sample from population; $\tau_i$ i.i.d with mean 0, variance $\sigma_\tau^2$.
    \item Errors $\varepsilon_{ij}$ i.i.d $N(0,\sigma^2)$ and independent of $\tau_i$.
\end{itemize}

Variance structure:
\begin{itemize}
    \item $\text{Var}(Y_{ij})=\sigma_\tau^2+\sigma^2$.
    \item $\text{Cov}(Y_{ij},Y_{ij'})=\sigma_\tau^2$ (same level). Different levels: covariance 0.
\end{itemize}

Expected mean squares:
\[
E[MS_E]=\sigma^2,\qquad E[MS_{\text{Trt}}]=\sigma^2+n\sigma_\tau^2.
\]

F-test remains $F=MS_{\text{Trt}}/MS_E$ to test $H_0:\sigma_\tau^2=0$.

\subsection{5.3 Variance Components \& ICC}

Estimators:
\begin{itemize}
    \item $\hat{\sigma}^2=MS_E$.
    \item $\hat{\sigma}_\tau^2=(MS_{\text{Trt}}-MS_E)/n$ (may be negative if true variance small; then set to 0 in practice).
\end{itemize}

Intraclass correlation:
\[
\rho=\frac{\sigma_\tau^2}{\sigma_\tau^2+\sigma^2}=\text{Corr}(Y_{ij},Y_{ij'}).
\]

Interpretation: fraction of total variance explained by group identity. If $\rho\approx0$, grouping has little effect; $\rho$ near 1 means within-group units nearly identical.

\subsection{5.4 Example: Sires \& Calf Weights}

5 sires (random), 8 calves each: $N=40$. From ANOVA: $MS_{\text{Trt}}=1397.8$, $MS_E=463.8$, $F=3.014$, p$\approx0.031$.

$\hat{\sigma}^2=463.8$, $\hat{\sigma}_\tau^2=(1397.8-463.8)/8=116.75$, ICC $=116.75/(116.75+463.8)=0.201$.

CI for overall mean (random): $82.55\pm t_{0.025,35}\sqrt{MS_{\text{Trt}}/N}\approx(70.5,94.6)$; fixed-effects CI using $MS_E$ would be narrower.

\begin{calcbox}
TI-84: \texttt{Fcdf(3.014,1E99,4,35)} for F p-value. For CI on $\sigma^2$, use $\chi^2$ distribution with df $=N-a$.
\end{calcbox}

%% ============================================================
\section{6. RCBD (Blocking)}
%% ============================================================

\subsection{6.1 Motivation, Model, Assumptions}

Blocking: group experimental units into homogeneous blocks for nuisance factors (location, day, operator) and randomize treatments within each block.

RCBD model (fixed blocks):
\[
y_{ij}=\mu+\tau_i+\beta_j+\varepsilon_{ij},\quad i=1,\ldots,a,\; j=1,\ldots,b.
\]

Assumptions:
\begin{itemize}
    \item Errors $\varepsilon_{ij}$ independent, $N(0,\sigma^2)$.
    \item Additive model: no treatment$\times$block interaction; difference between any two treatments same across all blocks.
\end{itemize}

\subsection{6.2 ANOVA Table \& df}

\begin{center}\small
\begin{tabular}{@{}lccc@{}}
\toprule
Source & df & SS & F \\
\midrule
Treatment & $a-1$ & $SS_{\text{Trt}}$ & $MS_{\text{Trt}}/MS_E$ \\
Block & $b-1$ & $SS_{\text{Block}}$ & $MS_{\text{Block}}/MS_E$ \\
Error & $(a-1)(b-1)$ & $SS_E$ & \\
Total & $ab-1$ & $SS_T$ & \\
\bottomrule
\end{tabular}
\end{center}

Error df: $(a-1)(b-1)$. For $a=4$, $b=3$, df$_E=(4-1)(3-1)=6$.

\subsection{6.3 Why Blocking Helps}

Apple orchard example (5 treatments, 6 locations):

\begin{center}\small
\begin{tabular}{@{}lcc@{}}
\toprule
Design & $MS_E$ & Conclusion \\
\midrule
CRD & 52,938 & No trt diff; high variance \\
RCBD & 34,455 & Same conclusion, but 35\% lower $MS_E$ \\
\bottomrule
\end{tabular}
\end{center}

Even if the overall conclusion is unchanged here, reduced $MS_E$ gives more power in other settings.

\subsection{6.4 Mixed RCBD: Random Blocks}

Random blocks model: $Y_{ij}=\mu+\tau_i+\beta_j+\varepsilon_{ij}$ with $\beta_j\sim N(0,\sigma_\beta^2)$ i.i.d, treatments fixed.

\begin{itemize}
    \item ANOVA table and dfs unchanged; treatment F-test still $MS_{\text{Trt}}/MS_E$.
    \item Block variance component: $\hat{\sigma}_\beta^2=(MS_{\text{Block}}-MS_E)/a$.
    \item Focus of inference: treatment means/effects; blocks model nuisance variability.
\end{itemize}

If interaction between blocks and treatments is real but not modeled, interaction gets absorbed into $SS_E$, inflating $MS_E$ and reducing power.

\subsection{6.5 Multiple Comparisons in RCBD}

Post-hoc tests for treatment means use same logic as CRD but with:
\begin{itemize}
    \item Error df $=(a-1)(b-1)$.
    \item Effective ``replicates'' per treatment $=b$ (one per block).
    \item $SE(\bar{y}_{i.}-\bar{y}_{j.})=\sqrt{2MS_E/b}$.
\end{itemize}

%% ============================================================
\section{7. Power \& Sample Size}
%% ============================================================

\subsection{7.1 Concepts}

Power: $P(\text{reject }H_0 \mid H_a\text{ true})$. Depends on $\alpha$, $n$, effect size, and $\sigma^2$.

\begin{itemize}
    \item Larger power (e.g., 0.9 vs 0.8) $\Rightarrow$ need larger $n$.
    \item Smaller $\alpha$ (e.g., 0.01 vs 0.05) $\Rightarrow$ larger $n$.
    \item Larger $\sigma^2$ $\Rightarrow$ larger $n$.
    \item Larger effect size $\Rightarrow$ smaller $n$.
\end{itemize}

\subsection{7.2 Cohen's $f$ for One-Way ANOVA}

Effect size:
\[
f=\sqrt{\frac{\frac{1}{a}\sum_{i=1}^a(\mu_i-\bar{\mu})^2}{\sigma^2}},\quad \bar{\mu}=\frac{1}{a}\sum\mu_i.
\]

Benchmarks: $f=0.10$ (small), $0.25$ (medium), $0.40$ (large). Required $n$ grows like $1/f^2$.

\subsection{7.3 Worked Example}

Means $(50,60,50,60)$, $a=4$, $\sigma^2=25$. $\bar{\mu}=55$. Between variance $=25$, so $f=1.0$ (very large). Standard power formulas give $n\approx3$ per group for 90\% power at $\alpha=0.05$.

If $\sigma=6$, $f=\sqrt{25/36}=0.833$; if $\sigma=7$, $f=\sqrt{25/49}=0.714$. As $\sigma$ doubles, $f$ halves and $n$ roughly quadruples.

\subsection{7.4 CI-Based Sample Size}

For CI on $\mu_i-\mu_j$ with half-width $ME$ at confidence $1-\alpha$:
\[
ME=t_{\alpha/2,df}\sqrt{2MS_E/n}.
\]

Solve for $n$, iterating because $df$ depends on $n$. For multiple intervals, may apply Bonferroni (replace $\alpha$ by $\alpha/K$).

%% ============================================================
\section{8. Non-Parametric Alternatives}
%% ============================================================

\subsection{8.1 When to Use}

Use when normality/variance assumptions badly violated and transformations fail. Trade lower power under normality for fewer assumptions.

\begin{center}\small
\begin{tabular}{@{}llll@{}}
\toprule
Parametric & Non-parametric & Key Assumption & Target \\
\midrule
Paired $t$ & Sign test & None & Median diff \(\neq0\) \\
Paired $t$ & Wilcoxon signed-rank & Symmetric diff dist & Location shift \\
Two-sample $t$ & Mann--Whitney & Similar shapes & Location shift / $P(X>Y)$ \\
ANOVA $F$ & Kruskal--Wallis & Continuous data & Equal distributions \\
\bottomrule
\end{tabular}
\end{center}

\subsection{8.2 Wilcoxon Signed-Rank}

For paired data or one-sample median test. Steps: compute $d_i=x_i-m_0$, drop zeros; rank $|d_i|$; sign ranks; sum positive ranks $W^+$. Under $H_0$, $E(W^+)=n(n+1)/4$, $\text{Var}(W^+)=n(n+1)(2n+1)/24$; use normal approx for moderate $n$.

Assumption: distribution of differences is symmetric; no normality needed.

\subsection{8.3 Mann--Whitney Rank-Sum}

Two independent groups sizes $m,n$. Rank pooled data; let $W_1$ be rank sum for group 1. Convert to $U_1=W_1-m(m+1)/2$. Under $H_0$, $E(U_1)=mn/2$, $\text{Var}(U_1)=mn(m+n+1)/12$.

Assumptions: independent samples, continuous distributions, similar shapes across groups; tests location shift.

\subsection{8.4 Kruskal--Wallis $H$}

Rank all $N$ observations; let $T_i$ be rank sum in group $i$ of size $n_i$.
\[
H=\frac{12}{N(N+1)}\sum_{i=1}^a\frac{T_i^2}{n_i}-3(N+1),
\]
approx $\chi^2_{a-1}$ under $H_0$.

Assumptions: independent samples, continuous distributions. Under true normality with equal variances, ANOVA $F$ more powerful; under heavy skew/outliers, K-W can be preferred.

\begin{calcbox}
TI-84: for $\chi^2$ p-value use \texttt{chi2cdf(H,1E99,a-1)}; for normal approximations use \texttt{normalcdf}.
\end{calcbox}

%% ============================================================
\section{9. Experimental Design Principles}
%% ============================================================

\subsection{9.1 Three Principles}

\begin{itemize}
    \item Randomization: randomly assign treatments to units; protects against bias, justifies independence.
    \item Replication: multiple units per treatment; estimates error variance and increases precision.
    \item Blocking: group similar units into blocks; removes block-to-block variability from error term.
\end{itemize}

\subsection{9.2 Design Comparison}

\begin{center}\small
\begin{tabular}{@{}p{1.7cm}p{2.3cm}p{3.0cm}p{2.3cm}@{}}
\toprule
Design & Nuisance factor & Model & Error df \\
\midrule
CRD & None & $y_{ij}=\mu+\tau_i+\varepsilon_{ij}$ & $N-a$ \\
RCBD & 1 block & $y_{ij}=\mu+\tau_i+\beta_j+\varepsilon_{ij}$ & $(a-1)(b-1)$ \\
Latin Sq & 2 blocks & $y_{ijk}=\mu+\tau_j+\rho_i+\gamma_k+\varepsilon$ & $(p-1)(p-2)$ \\
\bottomrule
\end{tabular}
\end{center}

\subsection{9.3 Quick Concept Checks}

\begin{itemize}
    \item ANOVA treats quantitative and qualitative factors the same in SS computations. (True)
    \item Fisher LSD is extremely conservative. (False; it is liberal.)
    \item Tukey HSD compares all treatments to a control. (False; Dunnett does.)
    \item $MS_E$ estimates the random error variance $\sigma^2$. (True.)
    \item In mixed RCBD with random blocks, inference focuses on treatments, not blocks. (True.)
    \item Power of Kruskal--Wallis is less than F-test under normality. (True.)
\end{itemize}

%% ============================================================
\section{10. Common Distributions \& TI-84 Reference}
%% ============================================================

\begin{center}\small
\begin{tabular}{@{}llll@{}}
\toprule
Distribution & Mean & Variance & Use \\
\midrule
$N(\mu,\sigma^2)$ & $\mu$ & $\sigma^2$ & Errors, CLT \\
$t_\nu$ & 0 & $\nu/(\nu-2)$ ($\nu>2$) & Small-sample means \\
$F_{d_1,d_2}$ & $d_2/(d_2-2)$ ($d_2>2$) & Variance ratios, ANOVA \\
$\chi^2_\nu$ & $\nu$ & $2\nu$ & Variance CIs, K-W \\
Poisson($\lambda$) & $\lambda$ & $\lambda$ & Counts; \(\sqrt{Y}\) transform \\
Binomial($n,p$) & $np$ & $np(1-p)$ & Proportions; arcsin$\sqrt{Y}$ \\
\bottomrule
\end{tabular}
\end{center}

\begin{calcbox}
TI-84 distributions (\texttt{2nd$\to$DISTR}): normal: \texttt{normalcdf}, \texttt{invNorm}; t: \texttt{tcdf}, \texttt{invT}; F: \texttt{Fcdf}; $\chi^2$: \texttt{chi2cdf}.\\
Example: upper-tail p-value for F=14.52, df=(2,21): \texttt{Fcdf(14.52,1E99,2,21)}.
\end{calcbox}

\vspace{2pt}\hrule\vspace{2pt}
\begin{center}\small
Use margins and blank areas for your own textbook examples and quick formulas.
\end{center}

\end{document}
